\documentclass[12pt]{article}
\usepackage[utf8]{inputenc}
\usepackage[T1]{fontenc}
\usepackage{hyperref}
\title{Waris Language Resources for K-12 Curriculum Development}
\author{Compiled by Claude Sonnet 5.5 (claude-sonnet-5-5) for Andr\'as Kornai}
\date{}
\begin{document}
\maketitle

\section{Introduction}
Waris (also called Walsa; ISO 639-3 wrs; Glottocode wari1266) is a Border-family language spoken in Sandaun Province, Papua New Guinea, and in Papua Province, Indonesia. Searches found only linguistic reference material. No dictionary, primer or school programme turned up.

\section{Primary Educational Resources}
No primers, textbooks or school materials were found in search results.

\section{Linguistic and Reference Materials}
\begin{itemize}
\item Wikipedia entry for Waris, which lists Robert Brown's papers (1981, on semantic aspects of Waris predications; 1988, on the case system and verb classification). I did not open the papers themselves. \url{https://en.wikipedia.org/wiki/Waris_language}
\item Glottolog entry: \url{https://glottolog.org/resource/languoid/id/wari1266}
\item Waris numeral data (Brown): \url{https://lingweb.eva.mpg.de/channumerals/Waris.htm}
\end{itemize}

\section{Religious and Community Resources}
A search summary suggested a Bible translation exists. I saw no URL for it, so it is not listed in the structured data.

\section{Digital and Audio Resources}
The Rosetta Panglossia page for wrs was returned in results. I did not check its content: \url{https://rosettapanglossia.longnow.org/index.php?title=Wrs}. Omniglot has a page on Waris writing: \url{https://omniglot.com/writing/waris.htm}.

\section{Curriculum Development Context}
The language is reported as stable, but no teaching materials were found. Development would need new documentation and orthography work with the community. Papua New Guinea and Indonesia would each have their own educational contexts.
\end{document}